\chapter{Declaraciones}
\label{cap:declaraciones}

Este cap\'itulo re\'une las nueve declaraciones obligatorias del informe,
cada una con secci\'on propia: contribuciones de autor\'ia, \'etica,
disponibilidad de datos, disponibilidad de c\'odigo, disponibilidad de
materiales, financiamiento, conflicto de intereses, uso de asistentes
autom\'aticos y agradecimientos.

\section{Contribuciones de autor\'ia (CRediT)}

La matriz completa de roles CRediT~\cite{credit-taxonomy}, verificada
individuo por individuo contra commits, ramas de trabajo y archivos
concretos del repositorio (\texttt{git shortlog -sne --all}), vive en
\texttt{CONTRIBUTORS.md}. La tabla~\ref{tab:credit-final} la resume:

\begin{table}[H]
  \centering
  \small
  \caption{Resumen de contribuciones CRediT por integrante (detalle completo en \texttt{CONTRIBUTORS.md}).}
  \label{tab:credit-final}
  \begin{tabular}{@{}p{3.6cm}p{4.4cm}p{5.4cm}@{}}
    \toprule
    Integrante & \'Areas principales & Evidencia \\
    \midrule
    Mariscal Cabrera, Jaime Josu\'e & Seguridad, autenticaci\'on/autorizaci\'on, cobertura JaCoCo, an\'alisis est\'atico/din\'amico de seguridad, ADR-006/007, CI & \rutalarga{Backend/src/main/java/com/biopet/security/}, \rutalarga{docs/mediciones/sec/}, \rutalarga{.github/workflows/} \\
    Beltr\'an Montiel, Fred Adri\'an & Procedimientos almacenados, acceso a datos, cach\'e, rendimiento k6, despliegue/GHCR & \rutalarga{db/procs/}, \rutalarga{docs/mediciones/perf/}, \rutalarga{docs/despliegue/} \\
    Taipe Mora, Zaida Melissa & Frontend Angular, accesibilidad, requisitos (SRS), SUS, Lighthouse, PRISMA & \rutalarga{frontend/src/app/}, \rutalarga{docs/requisitos/}, \rutalarga{docs/mediciones/sus/} \\
    \bottomrule
  \end{tabular}
\end{table}

Los tres autores de BIOPET, con afiliaci\'on institucional \'unica, son:
Fred Adri\'an Beltr\'an Montiel (\texttt{fbeltranm@uteq.edu.ec}), Jaime
Josu\'e Mariscal Cabrera (\texttt{jmariscalc@uteq.edu.ec}) y Zaida Melissa
Taipe Mora (\texttt{ztaipem@uteq.edu.ec}), todos de la Universidad
T\'ecnica Estatal de Quevedo. Otros colaboradores hist\'oricos que
participaron en fases anteriores del programa acad\'emico no son autores de
esta entrega y no se listan como tales en \texttt{CITATION.cff} ni en este
informe.

\section{\'Etica}

BIOPET utiliza tres categor\'ias de datos: datos sint\'eticos de
desarrollo/prueba; datos emp\'iricos de evaluaci\'on generados por
herramientas automatizadas contra el propio sistema
(\texttt{docs/etica/ETHICS.md}); y los 18 registros anonimizados
P01--P18 de la evaluaci\'on de usabilidad SUS, que seg\'un declaraci\'on
del equipo corresponden a participantes reales (detalle a continuaci\'on).
El sistema no
almacena datos cl\'inicos reales de mascotas ni ha sido operado con datos
de clientes reales de una cl\'inica veterinaria. La prueba de usabilidad
SUS se aplic\'o sobre 18 registros (P01--P18), matem\'aticamente
reproducibles desde sus respuestas Q1--Q10
(\texttt{scripts/analisis-sus.py}); seg\'un declaraci\'on del equipo,
corresponden a participantes reales. El protocolo declarado
(\texttt{docs/etica/ETHICS.md}) exig\'ia consentimiento informado firmado
por cada participante (plantilla en
\texttt{docs/etica/consentimientos/plantilla.md}), pero una auditor\'ia
documental posterior a esta entrega constat\'o que esos formularios
individuales firmados no est\'an actualmente disponibles como evidencia
verificable. Esta situaci\'on se document\'o, sin fabricar evidencia
retrospectiva, mediante una constancia de regularizaci\'on firmada por
los tres integrantes del proyecto
(\texttt{docs/etica/regularizacion-sus/CONSTANCIA-REGULARIZACION-SUS-BIOPET-2026-08-31.pdf}),
que \textbf{no sustituye} consentimientos individuales ni incorpora
firmas de participantes. Los resultados agregados se identifican
\'unicamente por c\'odigo de participante (P01--P18). Ning\'un documento
de este informe ni sus figuras
reproducen contrase\~nas con valor, JWT completos, valores de cookies,
encabezados \texttt{Authorization} con valor real, claves privadas ni el
identificador \texttt{jti} de un token real.

\section{Disponibilidad de datos}

El conjunto de datos de evidencias emp\'iricas (rendimiento, seguridad,
usabilidad, calidad web, cobertura) est\'a archivado permanentemente en
Zenodo con DOI \texttt{\doidataset}, y disponible tambi\'en en
\rutalarga{docs/mediciones/} del repositorio, con procedencia documentada en
\rutalarga{docs/mediciones/DATA-PROVENANCE.md} y diccionario de variables en
\rutalarga{docs/mediciones/DATA-DICTIONARY.md}, siguiendo los principios
FAIR (\rutalarga{docs/checklists/fair.md}).

\section{Disponibilidad de c\'odigo}

El c\'odigo fuente completo (backend, frontend, migraciones, scripts) est\'a
disponible en \texttt{https://github.com/Grinjoww/ENTREGA-FINAL-BIOPET}
bajo licencia MIT, y archivado permanentemente en Zenodo con DOI
\texttt{\doisoftware}. La imagen de contenedor del backend est\'a publicada
en GHCR con referencia inmutable por \emph{digest}
(cap\'itulo~\ref{cap:despliegue}).

\textbf{Reproducibilidad.} Todos los resultados cuantitativos de este
informe son reproducibles con los comandos documentados en el
cap\'itulo~\ref{cap:despliegue} y en \texttt{docs/informe/README.md},
contra el mismo commit citado en cada reporte de evidencia
(\texttt{docs/mediciones/*/README.md} y \texttt{*.md} de reporte
individuales).

\section{Disponibilidad de materiales}

Adem\'as del c\'odigo y los datos de medici\'on (secciones anteriores),
este trabajo produjo materiales complementarios --- protocolos,
instrumentos de evaluaci\'on y configuraciones de auditor\'ia --- que
tambi\'en quedan versionados en el mismo repositorio y archivo Zenodo
citados arriba (DOI \texttt{\doisoftware}), sin un identificador aparte
porque forman parte del mismo commit archivado:

\begin{itemize}
  \item \textbf{Instrumento y protocolo SUS}: cuestionario de Brooke
    (1996)~\cite{sus-brooke1996} sin modificar, plantilla de
    consentimiento informado (\rutalarga{docs/etica/consentimientos/plantilla.md})
    y protocolo \'etico (\rutalarga{docs/etica/ETHICS.md}).
  \item \textbf{Configuraciones de auditor\'ia de calidad web}:
    \texttt{lighthouserc.js}, \texttt{lighthouserc.desktop.js},
    \texttt{lighthouserc.render.js} y \texttt{lighthouserc.render.desktop.js}
    (perfiles m\'ovil/escritorio, contra \texttt{localhost} y contra el
    despliegue p\'ublico de Render).
  \item \textbf{Modelo de arquitectura C4}: fuente Structurizr DSL
    versionada en \rutalarga{docs/diagrams/workspace.dsl}.
  \item \textbf{Decisiones de arquitectura (ADR)} en \rutalarga{docs/adr/}
    y especificaci\'on de despliegue reproducible
    (\texttt{docker-compose.yml}, \texttt{render.yaml}, \texttt{Makefile}).
  \item \textbf{Requisitos}: SRS versionado en \rutalarga{docs/requisitos/SRS.md}.
\end{itemize}

Ninguno de estos materiales incluye credenciales, secretos ni datos
personales identificables de participantes (ver secci\'on de \'Etica).

\section{Financiamiento}

Este proyecto es exclusivamente acad\'emico, desarrollado como Entrega
Final de la asignatura \asignatura{} en la Universidad T\'ecnica Estatal
de Quevedo. No recibi\'o financiamiento externo de ninguna agencia,
empresa ni entidad de fomento a la investigaci\'on (\texttt{CONTRIBUTORS.md},
categor\'ia CRediT \emph{Funding acquisition}: no aplica para ning\'un
integrante).

\section{Conflicto de intereses}

Ninguno de los tres autores declara conflicto de intereses en relaci\'on
con este trabajo.

\section{Uso de asistentes autom\'aticos}

El uso de herramientas de Inteligencia Artificial durante el desarrollo de 
este proyecto se limitó exclusivamente a funciones de apoyo y orientaci\'on acad\'emica. 
Estas herramientas fueron utilizadas como gu\'ia para la organizaci\'on y 
distribuci\'on de tareas entre los integrantes del equipo, as\'i como para la 
revisi\'on de avances y la identificaci\'on de posibles mejoras durante el proceso
 de desarrollo. La Inteligencia Artificial no sustituy\'o el trabajo de los 
 integrantes ni la toma de decisiones del equipo, siendo los autores responsables 
 del dise\~no, implementaci\'on, validaci\'on y resultados finales del proyecto.


\section{Agradecimientos}

El equipo agradece al docente de la asignatura \asignatura{},
\docente{} (\correodocente{}), por la gu\'ia acad\'emica y la
retroalimentaci\'on de evaluaci\'on que orientaron las correcciones de
esta entrega. Se agradece tambi\'en a la \universidad{} y a la
\facultad{} por el acceso a la infraestructura acad\'emica utilizada
durante el desarrollo del proyecto. No se recibi\'o apoyo, patrocinio ni
colaboraci\'on de ninguna persona, instituci\'on o empresa externa al
equipo de tres autores y al docente aqu\'i mencionados.
